\section{Discussion and future work}
\label{sec:limitations}
Our approach shares limitations with other methods that test for unobserved confounding. Since we rely on the transportability assumption, our test could misidentify violations of this assumption as unobserved confounding. 
In addition, the lower bound we provide is optimistic; outside the common support of the two studies, the unobserved confounding could be arbitrarily high. Furthermore, our test is designed to detect confounding structures that bias the average treatment effect and, hence, would not detect confounding bias that cancels out on average.



Our discussion suggests several important directions for future research. 
First, 
developing a more refined sensitivity model that accounts for the correlation between outcomes and unobserved confounders could result in a more powerful test. Second, our test could be adapted to the scenario where multiple observational datasets may be available but no randomized control trials. Lastly, it would be highly valuable to propose a procedure that not only identifies hidden confounding but also suggests specific interventions to mitigate it. 

\subsection*{Acknowledgements}
PDB was supported by 
the Hasler Foundation grant number 21050. JA was supported by the ETH AI Center. KD was supported by the ETH AI Center and the ETH Foundations of Data Science.

